\begin{flushright}
{\large ПРИЛОЖЕНИЕ 2}\\[0.5em]
Шаблон дневника практики
\end{flushright}

\vspace{1em}

\begin{center}
\textbf{\MakeUppercase{\UniversityName}}
\end{center}

\noindent\textbf{Направление:} \DirectionName\\
\textbf{ОП:} \ProgramName

\vspace{1em}

\begin{center}
{\Large\bfseries ДНЕВНИК НАУЧНО-ИССЛЕДОВАТЕЛЬСКОЙ ПРАКТИКИ}
\end{center}

\noindent за период с \PracticeStart{} по \PracticeEnd

\noindent\textbf{Студент:} \StudentName \hfill \textbf{номер группы:} \StudentGroup

\noindent\textbf{Руководитель практики:} \SupervisorName, \SupervisorRole

\vspace{1em}

\scriptsize
\setlength{\tabcolsep}{2pt}
\begin{longtable}{|L{1.9cm}|L{2.5cm}|L{3.7cm}|L{1.8cm}|L{2.3cm}|L{1.5cm}|}
\textbf{Дата} & \textbf{\shortstack[l]{Структурное\\ подразделение}} & \textbf{Краткое содержание работы} & \textbf{Вопросы} & \textbf{Результаты} & \textbf{\shortstack[l]{Отметка о\\ выполнении}} \\
\hline
11.02.2026 -- 20.02.2026 & СПбГУ, ПМ-ПУ & Уточнены тема практики и рамка проекта как системы поддержки принятия решений для выбора следующей культуры, согласована структура итогового отчета. & Как корректно ограничить выводы рамкой прототипа, а не оптимизатора. & Зафиксирована постановка задачи и перечень исследовательских этапов практики. & Выполнено \\
\hline
23.02.2026 -- 27.02.2026 & СПбГУ, ПМ-ПУ & Проанализированы этапы подготовки данных: очистка USDA NASS CSB/CDL, сопоставление культур и построение оконной выборки \texttt{history\_1, history\_2, history\_3 -> target}. & Какие классы считать нерелевантными и как сохранить интерпретируемое пространство целей. & Подтвержден подготовленный датасет на \PreparedRows{} строк и оконный датасет на \WindowRows{} строк до фильтрации редких классов. & Выполнено \\
\hline
02.03.2026 -- 10.03.2026 & СПбГУ, ПМ-ПУ & Проверены фильтрация редких целевых классов, итоговое пространство из \TargetClasses{} культур и фиксированное разбиение по \texttt{CSBID}. Выполнено сравнение простых baseline-подходов. & Как исключить утечку между обучающей, валидационной и тестовой выборками при наличии нескольких окон одного поля. & Зафиксирована итоговая выборка на \FilteredRows{} строк и воспроизводимое group-aware разбиение. & Выполнено \\
\hline
11.03.2026 -- 23.03.2026 & СПбГУ, ПМ-ПУ & Изучена финальная baseline CatBoost-модель: состав признаков, метрики качества и причины выбора в качестве основного предиктивного ядра. & Достаточен ли компактный набор признаков для устойчивого качества без усложнения проекта. & Подтверждено, что \FinalModel{} показывает на test Accuracy \BaselineAcc{}, Macro F1 \BaselineMacroFOne{} и Weighted F1 \BaselineWeightedFOne{}. & Выполнено \\
\hline
24.03.2026 -- 03.04.2026 & СПбГУ, ПМ-ПУ & Проанализированы промежуточные эксперименты: ветка с дополнительными признаками и гибрид \texttt{CatBoost + transition graph}. & Может ли графовый сигнал улучшить ранжирование без потери качества по сравнению с baseline CatBoost. & Зафиксировано, что лучший гибрид с $\alpha=\GraphAlpha$, $\beta=\GraphBeta$ немного уступает baseline и не заменяет финальную модель. & Выполнено \\
\hline
06.04.2026 -- 14.04.2026 & СПбГУ, ПМ-ПУ & Подготовлены итоговые материалы производственной практики: краткий отчет, дневник и артефакты работы. & - & Сформирован комплект итоговых материалов. & Выполнено \\
\hline
\end{longtable}
\setlength{\tabcolsep}{6pt}
\normalsize

\vspace{0.5em}

\noindent * Подпись руководителя практики от организации
